%!TEX root = ../main.tex
% =============================================================================
%  三、数据预处理与特征分析      建议 2–3 页
%  评审在问：数据你读懂了吗？
%  本节必须出现：① 时标与对齐结论  ② 附件 1 身份核查（= 全年逐时段均值，非预报）
%                ③ 相关性的尺度拆解  ④ 预报误差分档（全时段 / 白天双口径）
%                ⑤ 数据质量与缺失值结论
%  禁止：相关性写成因果（“导致”“引起”“造成”一律换成“共动”“同向变动”）
% =============================================================================

\section{数据预处理与特征分析}

\subsection{数据来源与基本规模}

四个附件的用途见表~\ref{tab:data}。时间序列均按 \textbf{10~min} 粒度、以\textbf{时段末端}
为标签对齐，全天 144 个时段；负载与光伏为\textbf{功率}（kW）、电价为\textbf{元/kWh}，
不得混用。



\subsection{时标语义与对齐}

附件 1、2、4 的时段标签为 \(0{:}10,0{:}20,\dots,23{:}50,0{:}00{+}1\)（后者为次日
\(0{:}00\)），本文采用\textbf{位置读法}：

\begin{itemize}
  \item 标签表示\textbf{时段末端}：槽位 \(j\)（\(j=0,1,\dots,143\)）表示区间
        \([10j,\ 10(j+1))\) 分钟；第 \(j+1\) 个数据行/列对应槽位 \(j\)，标签只作显示、
        \textbf{不参与重映射}；
  \item 每行 144 个值\textbf{无缝铺满} \(0{:}00\) 至 \(24{:}00\)；末列 \(0{:}00{+}1\) 属当日
        \(23{:}50\)--\(24{:}00\)、\textbf{不得丢弃}，否则合计偏低 \(1/144\approx 0.6944\%\)。
\end{itemize}

附件 3 的列标签为“预报 1 小时”至“预报 24 小时”，约定\textbf{预报 \(k\) 小时对应发布
时刻之后的第 \(k\) 个小时}，即 \([h_0+k-1,\ h_0+k)\)（\(h_0\in\{0,6,12,18\}\)）；若误按
\([h_0+k,\ h_0+k+1)\) 对齐，则\textbf{全部预报整体错 1 小时}。\(18{:}00\) 的预报中
\(k=7,\dots,24\) 落在\textbf{次日} \(0{:}00\)--\(17{:}00\)（18 个整点），跨日用绝对时间轴
映射（目标整点序号 \(d\times 24+h_0+k-1\)），不丢配对样本。表~\ref{tab:timeruler} 给出题面
时段与附件标签的对照。

\begin{table}[htbp]
  \centering
  \caption{题面表 1 的时段写法与附件时段标签对照（示例；时段序号 \(i\) 从 1 起计，
           槽位下标 \(j=i-1\)）}
  \label{tab:timeruler}
  \small
  \begin{tabular}{cccc}
    \toprule
    题面写法 & 附件标签 & 对应区间 & 数据点序号 \(i\) \\
    \midrule
    \(10{:}00\)--\(10{:}10\) & \(10{:}10\) & \([10{:}00,10{:}10)\) & 61 \\
    \(12{:}00\)--\(12{:}10\) & \(12{:}10\) & \([12{:}00,12{:}10)\) & 73 \\
    \(14{:}00\)--\(14{:}10\) & \(14{:}10\) & \([14{:}00,14{:}10)\) & 85 \\
    \(16{:}00\)--\(16{:}10\) & \(16{:}10\) & \([16{:}00,16{:}10)\) & 97 \\
    \(20{:}00\)--\(20{:}10\) & \(20{:}10\) & \([20{:}00,20{:}10)\) & 121 \\
    \(23{:}50\)--\(24{:}00\) & \(0{:}00{+}1\) & \([23{:}50,24{:}00)\) & 144 \\
    \bottomrule
  \end{tabular}
\end{table}

对齐结果由两条独立校验确认：结构断言（附件 1 恰 144 行；附件 3 前向填充后每天恰 4 行）与结果
文件层面的汇总校验——审计项 A10 核对“全天购电量/全天购电费”与逐时段求和，实测最大相对偏差为
\(3.40\times10^{-7}\)（问题 2、3）与 \(4.32\times10^{-7}\)（问题 4-2），远低于
\(10^{-4}\) 的阈值。该汇总校验对“整体错位一个槽”不可见，该风险由位置读法与
表~\ref{tab:timeruler} 逐格对照锁定。

\subsection{附件 1 的身份核查}

附件 1 第 4 列的表头写作“光伏发电\textbf{预测}功率”，但逐项核查表明它并非预报：

其\textbf{光伏}、\textbf{负载}、\textbf{电价}列分别等于附件 2 全年（365 天）逐时段实际光伏功率、
负载与附件 4 逐时段电价的均值，最大偏差分别为 \(0.0431\)~kW、\(5.04\times10^{-5}\)~kW 与
\(5.15\times10^{-5}\)~元/kWh。

三处偏差都在数值精度量级（\(10^{-5}\) 及以下），故附件 1 是\textbf{全年平均日（典型日）}而非
日前预报（电价列的一致同时是假设 A10“分时骨架”的数据依据）：附件 1 \textbf{只能}用于问题 1
与基线口径，\textbf{不可}用作问题 2--4 的预报输入或验证集；问题 1 因此是
\textbf{确定性优化问题}，其最优值是问题 2--4 的对照基准。

换用窗口口径（2025-02-01 至 12-31 共 334 天）的逐时段均值作比较，电价偏差升到
\(3.07\times10^{-2}\)~元/kWh，为\textbf{全年 365 天}口径的
\(3.07\times10^{-2}/5.15\times10^{-5}\approx600\) 倍：\textbf{全年口径}支撑“附件 1 即
附件 2/4 逐时段均值”的身份断言，\textbf{窗口口径}则说明问题 1 的典型日结果\textbf{不可}
直接当作问题 2--4 的窗口基准（§1.3 表~\ref{tab:errorprop} 第二行）。

\subsection{电价与负载的结构特征}

按附件 4 统计，波动电价均值 \(0.7662\)~元/kWh、最小 \(0.0076\)、最大 \(1.7936\)~元/kWh；
按附件 1 统计，固定电价均值同为 \(0.7662\)~元/kWh、最小 \(0.3713\)、最大 \(1.3952\)~元/kWh。
按附件 2 统计，负载均值 \(4626.0337\)~kW、最小 \(1995.7176\)、最大 \(7978.8849\)~kW，日均
用电量 \(110\,802.6332\)~kWh；光伏均值 \(2311.7857\)~kW、峰值 \(10\,216.2000\)~kW，日均
发电量 \(57\,092.4850\)~kWh。

两类结构性事实直接影响建模方式。其一，\textbf{套利是否有利可图}：盈亏平衡价比阈值为
\(1/(\eta_c\eta_d)=1/0.81\approx1.2346\)，而固定电价的极值之比
\(1.3952/0.3713\approx3.76\)（口径见 §1.1）远高于它。其二，\textbf{约束落在哪一侧}：光伏峰值
\(10\,216.2000\)~kW 明显高于充放电功率上限 \(5000\)~kW，故有光伏功率超出充电能力的时段，
盈余只能靠储能吸纳或弃光；容量摆幅与功率谁更紧由第四章的容量/功率双向扫描判定。

电价与负载的关系具有\textbf{尺度依赖性}（表~\ref{tab:corr} 按三个尺度报告）。

\begin{table}[htbp]
  \centering
  \caption{电价与负载、净负荷、光伏的相关系数（三个尺度；池化与去形状口径的样本量为
           \(365\times144=52\,560\) 个时段观测，日级口径为 365 个日观测；
           全样本为 2025 年全年 365 天口径，
           与情景保真度检验（\(0.2620\rightarrow0.2597\)）验证，见 §7.2）}
  \label{tab:corr}
  \small
  \setlength{\tabcolsep}{4pt}
  \begin{tabular}{lccc}
    \toprule
    变量对 & 池化（全部时段） & 日级（逐日先算再合并） & 剔除日内形状后 \\
    \midrule
    电价 --- 负载 & \(+0.5077\) & \(+0.8977\) & \(+0.4554\) \\
    电价 --- 净负荷 \(L-PV\) & \(+0.2602\) & \(+0.9824\) & \(+0.1583\) \\
    电价 --- 光伏 & \(+0.0284\) & \(-0.3431\) & \(+0.0484\) \\
    \bottomrule
  \end{tabular}
\end{table}

读法如下。① \textbf{电价与负载共动}：日级 \(+0.8977\)，剔除日内共同形状后仍保留
\(+0.4554\)。② \textbf{“电价—光伏”的池化值具有误导性}：池化仅 \(+0.0284\)、日级
\(-0.3431\)，方向相反；真实结构是\textbf{光伏出力大的午间对应电价低谷}，独立抽样会生成
“午间光伏大发同时电价高企”的不可能情景并低估尾部费用，故问题 4
用保留相关结构的联合情景（§7.2）。③ “电价—净负荷”随尺度递进：池化 \(+0.2602\)、日级
\(+0.9824\)、剔除日内形状后 \(+0.1583\)；④ “电价—光伏”去形状后为 \(+0.0484\)，接近 0，说明
同向关系几乎全部来自日尺度。以上只描述\textbf{共动/同向变动}，本文不主张其中任何一方为因。

\subsection{光伏预报误差的统计刻画}

本节把附件 3 中\textbf{当天 \(0{:}00\) 发布}的 24~h 预报（即问题 2/4-2 的日前计划依据）与
附件 2 实际功率逐槽配对：整点预报先展开为 6 个 10~min 槽，再逐槽相减得
\(r=\widehat{PV}-PV\)（单位 kW）。夜间光伏接近 0、相对误差无意义，故同时报告两个口径，并统一
以\textbf{目标整点} \(h\ge5\)（\(5{:}00\)--\(24{:}00\)，每天 19 个整点）为白天口径、不按功率
阈值筛选，使两口径分母固定、可比（表~\ref{tab:fcsterr}）。

\begin{table}[htbp]
  \centering
  \caption{光伏预报误差统计（样本为 2025 年每天 \(0{:}00\) 发布的 24~h 预报与附件 2 对应
           实际功率的逐槽配对；MAE、RMSE、bias 单位均为 kW）}
  \label{tab:fcsterr}
  \small
  \begin{tabular}{lcccc}
    \toprule
    口径 & 样本量 \(n\) & MAE & RMSE & bias \\
    \midrule
    全时段（目标整点 \(0\)--\(23\) 时，144 槽/天） & \(52\,560\) & \(361.58\) & \(636.37\) & \(-26.96\) \\
    白天（目标整点 \(\ge5\) 时，114 槽/天） & \(41\,610\) & \(456.72\) & \(715.22\) & \(-34.06\) \\
    \bottomrule
  \end{tabular}
\end{table}

白天口径的 MAE 为全时段口径的 \(456.72/361.58\approx1.26\) 倍（夜间接近 0 的样本摊薄了
全时段指标），建模应以白天口径为依据。更关键的是\emph{bias 的符号}：按目标整点聚合后正负号
系统性变号，见表~\ref{tab:biasprofile}。

表~\ref{tab:biasprofile} 的形状有直接的决策含义：\textbf{上午被系统性高估}（\(7\) 时达
\(+733.00\)~kW）使 \(0{:}00\) 计划的光伏出力偏大、购电量偏少，缺口更易触发 \(5\) 倍紧急购电
（口径 \calref{C-05}）；\textbf{下午被系统性低估}（\(15\) 时达 \(-801.87\)~kW）使计划购电偏多，
多买电量按计划电价白付，或在问题 3 中触发 \(0.5\) 倍违约电价与 \(1.5\) 倍超额电价。全天聚合
bias 只有 \(-26.96\)~kW（白天 \(-34.06\)~kW），是两类相反偏差抵消后的净额，\textbf{不能据此
判断风险方向}；\(19\)--\(23\) 时的平均偏差落在 \([-0.15,0.00]\)~kW 内，出力均接近 0，
不触发任何惩罚。故决策层必须按\textbf{目标整点}设置安全边际而非按全天平均误差一次性纠偏，
场景生成直接用逐槽预报残差的经验分布（§5.2）以保留这种双向结构（\evid{C-Q2-01}）。

\paragraph{数据质量与缺失值}
三条结构断言全部通过：附件 1 恰 144 行；附件 3 日期列前向填充后每天恰有 4 个发布时刻
（合计 1460 行）；结果窗口恰 334 天（2025-02-01 至 12-31），结果表恰 147 列（日期 \(+\) 144
个时段 \(+\) 全天购电量 \(+\) 全天购电费）。断言通过后不再插补：附件 3 稀疏日期列以
\texttt{ffill} 补齐，其余附件按位置读法逐列读入。数据侧唯一实质性质量问题是\textbf{预报质量的
非均匀性}（上午与下午偏差方向相反、傍晚之后无有效预报信息、误差水平白天高于全时段），故本文
不采用“预报误差独立同分布且无偏”的简化，而把逐槽残差的经验分布直接用于场景生成
（§5.2、§7.2），并在决策层按目标整点设安全边际与滚动更新（§5.4、§6.3）。
